
% --- WAT language definition and style ---
\lstdefinelanguage{WAT}{
  morekeywords={
    module,type,func,param,result,export,import,local,global,data,
    i32,i64,f32,f64,v128,const,set,get,load,store,unreachable,
    block,loop,if,else,br,br_if,call,offset,align,externref,ref,null
  },
  sensitive=true,
  alsoletter={.,_},
  comment=[l]{;;}
}

\lstdefinestyle{watpaper}{
  language=WAT,
  basicstyle=\ttfamily\scriptsize,
  keywordstyle=\color{blue!80!black}\bfseries,
  commentstyle=\itshape\color{green!50!black},
  stringstyle=\color{red!70!black},
  numbers=left,
  numberstyle=\tiny\color{black},
  stepnumber=1,
  numbersep=2pt,
  frame=lines,
  rulecolor=\color{black},
  framesep=6pt,
  xleftmargin=0pt,
  framexleftmargin=0pt,
  showstringspaces=false,
  keepspaces=true,
  columns=fullflexible,
  breaklines=true,
  breakatwhitespace=false,
  escapeinside={(*@}{@*)},
}

% A generic code style (same look as watpaper) usable with any language via [language=...]
\lstdefinestyle{codepaper}{
  basicstyle=\ttfamily\scriptsize,
  keywordstyle=\color{blue!80!black}\bfseries,
  commentstyle=\itshape\color{green!50!black},
  stringstyle=\color{red!70!black},
  numbers=left,
  numberstyle=\tiny\color{black},
  stepnumber=1,
  numbersep=2pt,
  frame=lines,
  rulecolor=\color{black},
  framesep=6pt,
  xleftmargin=0pt,
  framexleftmargin=0pt,
  showstringspaces=false,
  keepspaces=true,
  columns=fullflexible,
  breaklines=true,
  breakatwhitespace=false,
  escapeinside={(*@}{@*)},
}
% --- helper macro: tight colorbox ---
% \newcommand{\tightbox}[2]{\begingroup\fboxsep=0pt\colorbox{#1}{#2}\endgroup}

\begin{wrapfigure}[27]{r}{0.46\columnwidth}
% \vspace{-3ex}
\centering
\lstset{style=codepaper}
\begin{lstlisting}[language=C]
// Factorial in C
long long factorial(long long n) {
  if (n == 0) {
    return 1;
  } else {
    return n * factorial(n - 1);
  }
}
\end{lstlisting}
\centering\footnotesize{(a) C code.}
\vspace{4pt}
\lstset{style=codepaper}
\begin{lstlisting}[language=WAT]
(module
  (type (func (param i64) (result i64))) ;; Type section
  (func (type 0) (param i64) (result i64)
    local.get 0
    i64.eqz
    if (result i64) (i64.const 1) else
      local.get 0
      local.get 0
      i64.const 1
      i64.sub
      call 0
      i64.mul
    end)
  (export "factorial" (func 0)) ;; Export section
)
\end{lstlisting}
\centering\footnotesize{(b) Corresponding WAT code.}
\caption{Factorial example.}
\label{ex:background_ex1}
\end{wrapfigure}

% \begin{figure}[t]
% \centering
% % ----------- TOP ROW (2 columns) -----------
% \begin{minipage}[t]{0.48\linewidth}
% \lstset{style=codepaper}
% \begin{lstlisting}[language=C]
% // Factorial in C


% long long factorial(long long n) {


%   if (n == 0) 
%   {
%     return 1;
%   } 
    
%   else 
%   {
  
%    return n * factorial(n - 1);

%   }

% }
% \end{lstlisting}
% \centering\footnotesize{(a) C code to compute Factorial}
% \end{minipage}\hfill
% \begin{minipage}[t]{0.48\linewidth}
% \lstset{style=codepaper}
% \begin{lstlisting}[language=WAT]
% (module
%   (type (func (param i64) (result i64))) ;; -> Type section
%   (func (type 0) (param i64) (result i64) ;; -> Function + Code sections
%     local.get 0
%     i64.eqz
%     if (result i64) (i64.const 1) else
%       local.get 0
%       local.get 0
%       i64.const 1
%       i64.sub
%       call 0
%       i64.mul
%     end)
%   (export "factorial" (func 0))           ;; -> Export section
% )
% \end{lstlisting}
% \centering\footnotesize{(b) Corresponding WAT code}
% \end{minipage}

% \vspace{6pt}
% \caption{\textbf{Example of the factorial computation in C (a) and in WebAssembly Text (WAT) format (b) \cite{wiki_wasm}}
% }
% \label{ex:background_ex1}
% \end{figure}
